\section{The actual delta to Forge}
\label{sec:delta}

\subsection{A pinned, deliberately narrow novelty claim}
The reference point is Forge commit
\code{c98e47c5f804e92880fc1d0e378c1b95832b685c}. Its README and merged
provenance ledger describe eighteen proposals and an extensive existing
capability set \cite{forge,forge-provenance}. This article was prepared after
reading that merged ledger, its conclusions, the Leant/Djex integration notes,
and the neighboring projects' published interfaces. Targeted repository
searches supplemented that inspection. It is not a byte-by-byte audit of every
file in every frozen submission, and the claim of nonduplication is relative to
the published merged design and capability ledger.

In particular, induction with generalization, finite-algebra covers, polynomial
and observable-space closure, exact arithmetic witnesses, Kripke countermodels,
telescoping, Ore transport, cyclic proof compilation through well-founded
measures, and the certificate acceptance boundary are already baseline
material. None is presented here as a new invention. The baseline's own
conclusion correctly asks for narrow implementations instead of further vague
architecture \cite{forge-conclusion}. This proposal follows that direction by
shipping three complete Python search/replay pairs and by stating the precise
Lean theorems still required.

The new mathematical ingredients are also not original discoveries in the
literature: pushdown saturation, finite B\"uchi games, and finite-chain
absorption are established subjects. What is new to the inspected Forge design
is a set of \emph{source-level obligations}, \emph{specialized certificate
formats}, and \emph{implemented adapters and tests} for these subjects. The
article makes no claim of worldwide priority for the algorithms or rank
arguments.

\begin{table}[ht]
\centering\small
\begin{tabularx}{\textwidth}{@{}p{0.23\textwidth}XX@{}}
\toprule
Existing prerequisite & Added capability & Distinction that must survive translation \\
\midrule
Finite-algebra closure & Pushdown reachability and regular stack targets & A finite summary describes infinitely many configurations; it is not a finite enumeration of those configurations. \\
Well-founded cyclic proof compilation & Recurrent adversarial strategies & Execution may continue forever. A rank resets after acceptance rather than decreasing globally. \\
Weighted finite-word observations & Infinite-horizon probabilistic values & A harmonic equation alone does not prove absorption, uniqueness, or preservation of probability mass. \\
Typed witness synthesis & Strategies with persistent behavioral laws & One policy must work against every future opponent choice, not one unrelated witness per finite sample. \\
\bottomrule
\end{tabularx}
\caption{The extension boundary. Baseline mechanisms can be reused as infrastructure; they are not claimed to be unable to encode these problems in principle.}
\label{tab:delta}
\end{table}

There is an important qualification to the first row. Pop summaries can be
viewed as a finite relational closure, and derivations can be encoded by trees.
Consequently Forge's generic cover or closure principles may be useful in their
formalization. The addition is not ``finite closure, again.'' It is the
suffix-parametric pushdown semantics, its exact translation, its productive
proof representation, its regular-target compiler, and the resulting complete
worker for a named infinite-state fragment. The same distinction applies to
reusing ordinary arithmetic to check ranks and rational residuals.

\subsection{Three semantic quantifiers rather than three solver names}
Consider a program that can remain in a loop or leave it. Depending on the
meaning of a transition, three entirely different questions arise:
\begin{align*}
&\exists\text{ a finite execution reaching a target};\\
&\exists\sigma\;\forall\text{ opponent choices}\;\GF\,B;\\
&\Prob(\text{termination})=1\quad\text{and}\quad\E[\text{cost}]=v.
\end{align*}
The first is existential reachability. The second asks for a persistent strategy
that visits $B$ infinitely often. The third is a statement about a measure on
executions. A solver that forgets which question it is answering can return a
perfectly valid certificate for the wrong proposition.

The proposed router therefore recognizes a \emph{behavioral contract}, not just
syntactic occurrences of recursion, a graph, or a matrix. An empty-stack query
can enter the pushdown lane. A universally quantified infinite play under an
existential policy can enter the game lane. A hitting-time or expected-cost
claim can enter the probabilistic lane only after transition probabilities and
terminal behavior have been preserved. Unsupported source constructs remain
outside these fragments; the worker must not manufacture a finite model by
silently erasing them.

\subsection{What is already executable}
The package uses the Python standard library only. All probabilities and
numerical certificates use \code{fractions.Fraction}; no floating-point solve
is used. There are three search modules and one replay module. Replay imports
neither search nor the algorithms used by the independent reference oracles.
The input dataclasses and rational arithmetic runtime are shared, so this is
algorithmic separation, not a proof of independent implementation correctness.

The payload of a successful Python replay is a \emph{research result}. It is
not a Lean expression, a theorem declaration, or a trusted receipt from the
existing Forge boundary. In the proposed Lean implementation, a successful
certificate check must still yield a proof of the corresponding object-level
semantic theorem and be transported back to the original goal. These are
explicitly separate milestones throughout the article.
